% ============================================================
% CAPITOLO 8 - DEPLOY E INTERFACCIA UTENTE
% ============================================================

\section{Messa in Produzione e Interfaccia Utente}

La fase conclusiva del progetto ``Signify'' ha riguardato l'ingegnerizzazione del modello in un'applicazione web moderna, fruibile in tempo reale direttamente dal browser senza la necessità di installare software aggiuntivo. 

\subsection{Architettura del Frontend (Vite + React)}
Il sistema di traduzione è stato integrato in una \textbf{Progressive Web App (PWA)} sviluppata con \textit{React} e gestita tramite il build tool \textit{Vite}, scelto per la sua estrema velocità di compilazione e gestione dei moduli. 

L'architettura del frontend si appoggia alle seguenti tecnologie:
\begin{itemize}
    \item \textbf{Mediapipe Holistic:} Utilizzato per l'estrazione dei landmark in tempo reale direttamente sul client (edge computing), riducendo la latenza e garantendo la privacy dell'utente.
    \item \textbf{ONNX Runtime Web:} Il modello LSTM addestrato in PyTorch è stato convertito nel formato standard \textit{ONNX} per essere eseguito con accelerazione hardware (WebGL/WebGPU) all'interno del browser.
    \item \textbf{Tailwind CSS:} Per la realizzazione di un'interfaccia utente pulita, accessibile e completamente reattiva.
\end{itemize}

\subsection{Responsive Design e Prova su Dispositivi}
L'applicazione è stata progettata seguendo il paradigma \textit{Mobile-First}, garantendo un'esperienza d'uso ottimale sia su desktop che su dispositivi mobili (Smartphone e Tablet). Questo è stato reso possibile grazie alla natura flessibile del CSS e all'ottimizzazione del carico computazionale dei modelli di Computer Vision.

Nelle figure seguenti è possibile osservare il comportamento dell'interfaccia su diverse risoluzioni, verificate tramite i \textit{Chrome DevTools}:

\begin{figure}[H]
    \centering
    \includegraphics[width=0.9\textwidth]{DevTools localhost 5173}
    \caption{Visualizzazione Desktop dell'interfaccia di traduzione. Il pannello laterale (HUD) mostra la probabilità del segno rilevato e la cronologia della frase.}
    \label{fig:deploy_desktop}
\end{figure}

\begin{figure}[H]
    \centering
    \begin{minipage}{0.48\textwidth}
        \centering
        \includegraphics[width=\textwidth]{DevTools localhost iPad Air (1)}
        \caption{Layout su iPad Air.}
        \label{fig:deploy_ipad}
    \end{minipage}
    \hfill
    \begin{minipage}{0.48\textwidth}
        \centering
        \includegraphics[width=\textwidth]{DevTools localhost iPhone 14 Pro Max}
        \caption{Layout su iPhone 14 Pro Max.}
        \label{fig:deploy_iphone}
    \end{minipage}
\end{figure}

\subsection{Strategia di Deploy}
Il deploy dell'applicativo è stato effettuato tramite \textbf{GitHub Pages}, sfruttando le \textit{GitHub Actions} per l'automazione del processo di build e distribuzione (CI/CD). Questa scelta garantisce una distribuzione globale tramite CDN e una manutenzione semplificata del codice sorgente.
